\documentclass{jsarticle}
\usepackage[dvipdfmx, final]{graphicx}
\usepackage{amsmath}
\usepackage{amssymb}
\usepackage{chemfig}
\begin{document}
\title{}
\author{}
\date{}
\maketitle
 $$  \chemfig{CH_3-CH(-[2]CH_3)-CH_2-C(=[:60]O)-[:-60]O-H} $$ 

 富山大学、1998年度の入学試験の有機化学の問題。
	$$ \chemfig{CH_3-C(=[:90]O)-0H} $$
	$CH_3 COOH$の分子量は、24+4+16=44であり、
	$\text{$二量体(88)\times n$}$、 N=4くらいだった。
	352くらいの分子量だった。
	$\text{$実測値=0.75\times 理論値=261mg$}$、
	収率=0.75だった。
	$\text{${xmg \over 264mg} < 1$}$、
	$\text{x=200}$くらいだった。
	264mgの化合物の構造式を導く問題だった。
	$\text{{$200 \over 264$} = 0.75}$、
	ある物質の構造式と、質量を知りたいので、
	原材料の264mgを、分析に掛けて見ると、酢酸が分子量が44mgであり、
	分析で無くなった質量が無いのと、収率の理論値と実存値の
	割合を求める問題である。その上に、誤差が規定外なので、　
	分子量を求めると、44mgであり、実材料は、261mgなのに、
	理論値が264mgになっていて、
	単体量では、誤差が大きすぎる、
	二量体にすると、
	$$ = {{44 \times 2 \times n} \over {264}} \le 1 $$
	$n = 4$くらいと、目星をつけた。$n=3$だった。
        収率は、標準では、1以下であることと、誤差が実存量から、
	水素結合になっていることから、
	二量体である。　
	収率を求めると、二量体が適切な解になっている。
	この問題は、理論値200mgの質量を求める問題だった。
	　
  2002年は、ノーベル化学賞に分析化学の田中耕一さんが、誤差のスペクトルを
	見抜くことを、この大学は、同じことをしていた。
\end{document}
